% !PRINTED 264
\weberpage{264}
Aus diesem System von Gleichungen leitet man ein zweites
und drittes her für die Funktionen $f_1$, $f_2$, indem man $\omega$ durch
$\omega + 1$ und darauf $\omega$ durch $-1:\omega$ ersetzt. Diese beiden Systeme
haben die gleiche Form, nur ist das eine Mal
\[
u_1 = f_1(\omega), \qquad v_1 = \left(\frac{2}{a}\right) f_1\left(\frac{c+\partial\omega}{a}\right),
\]
das andere Mal
\[
u_1 = f_2(\omega), \qquad v_1 = \left(\frac{2}{\partial}\right) f_2\left(\frac{c+\partial\omega}{a}\right)
\]
zu setzen. Aus den Vertauschungen (4) ergibt sich so:
{\small
\[
\begin{aligned}
\text{II.}\quad n &= 3, & A_1 &= \left(\frac{u_1}{v_1}\right)^6 - \left(\frac{v_1}{u_1}\right)^6, & B_1 &= (u_1v_1)^3 + \frac{8}{(u_1v_1)^3}, \\
& & A_1 + B_1 &= 0, \\
n &= 5, & A_1 &= \left(\frac{u_1}{v_1}\right)^3 - \left(\frac{v_1}{u_1}\right)^3, & B_1 &= (u_1v_1)^2 + \frac{4}{(u_1v_1)^2}, \\
& & A_1 + B_1 &= 0, \\
n &= 7, & A_1 &= \left(\frac{u_1}{v_1}\right)^4 + \left(\frac{v_1}{u_1}\right)^4, & B_1 &= (u_1v_1)^3 + \frac{8}{(u_1v_1)^3}, \\
& & A_1 - B_1 - 7 &= 0, \\
n &= 11, & A_1 &= \left(\frac{u_1}{v_1}\right)^6 - \left(\frac{v_1}{u_1}\right)^6, & B_1 &= u_1v_1 + \frac{2}{u_1v_1}, \\
& & A_1 + B_1^5 + B_1^3 - 2B_1 &= 0, \\
n &= 13, & A_1 &= \frac{u_1}{v_1} - \frac{v_1}{u_1}, & B_1 &= (u_1v_1)^6 + \frac{64}{(u_1v_1)^6}, \\
& & A_1^7 - 6A_1^5 + A_1^3 + 20A_1 + B_1 &= 0, \\
n &= 17, & A_1 &= \left(\frac{u_1}{v_1}\right)^3 + \left(\frac{v_1}{u_1}\right)^3, & B_1 &= (u_1v_1)^4 + \frac{16}{(u_1v_1)^4}, \\
& & A_1^3 - B_1^2 - 17A_1B_1 - 34A_1^2 - 34B_1 + 116A_1 \\
& & + 440 &= 0, \\
n &= 19, & A_1 &= \left(\frac{u_1}{v_1}\right)^2 - \left(\frac{v_1}{u_1}\right)^2, & B_1 &= (u_1v_1)^3 + \frac{8}{(u_1v_1)^3}, \\
& & A_1^5 + B_1^3 - 19A_1B_1^2 + 95A_1^2B_1 - 109A_1^3 \\
& & + 128B_1 - 128A_1 &= 0.
\end{aligned}
\]
}
